\section{Study Design}
\label{sec:design}

\subsection{Research Questions and Analysis Overview}

We organize the study around the following research questions:
\begin{description}
\item[RQ1: Outcomes.] To what extent do LLM agents reproduce known IoT
vulnerabilities under the studied conditions?
\item[RQ2: Breakdowns.] Where does reproduction break down, and what
observable conditions are associated with these breakdowns?
\item[RQ3: Responses.] How do agents respond to reproduction obstacles,
and how do their subsequent actions and reports relate to observed progress?
\end{description}

The analysis is exploratory: this framework was formulated while
reexamining completed experiments and preliminary records. It is not a
preregistered set of hypotheses. Figure~\ref{fig:study} separates execution
from the subsequent evidence analysis. RQ1 establishes what was observed;
RQ2 examines unmet conditions in those attempts; RQ3 studies subsequent
actions and their connection to artifacts and reports. None requires
access to a model's hidden internal reasoning.

\begin{figure*}[t]
\centering
% Adapted from the data / agent / evaluation organization of framework0724.pdf.
% The historical scoring panel is replaced by the current three analyses.
\begin{tikzpicture}[>=Latex,
  panel/.style={draw,rounded corners,align=left,inner sep=7pt,
    text width=1.93in,minimum height=1.05in,fill=black!3},
  analysis/.style={draw,rounded corners,align=center,inner sep=6pt,
    text width=1.93in,minimum height=0.65in}]
\node[panel,anchor=north west] (cases) at (0,0) {
  \textbf{Case construction}\\
  Public CVE sources\\
  30 selected cases\\
  Reference profiles};
\node[panel,anchor=north west] (agent) at (2.35in,0) {
  \textbf{Agent reproduction}\\
  CVE + common prompt\\
  Plan, execute, report\\
  Five configurations; three trials};
\node[panel,anchor=north west] (records) at (4.70in,0) {
  \textbf{Retained run evidence}\\
  Plans and reports\\
  Firmware and test artifacts\\
  Trajectories and service logs};
\draw[->] (cases.east)--(agent.west);
\draw[->] (agent.east)--(records.west);
\node[analysis,anchor=north west] (rq1) at (0,-1.56in) {
  \textbf{RQ1: Outcomes}\\Target identity, execution, effect};
\node[analysis,anchor=north west] (rq2) at (2.35in,-1.56in) {
  \textbf{RQ2: Breakdowns}\\Activities and unmet conditions};
\node[analysis,anchor=north west] (rq3) at (4.70in,-1.56in) {
  \textbf{RQ3: Responses}\\Actions, progress, and reports};
\draw[->] (records.south)--(rq3.north);
\draw[->] (records.south) -- ++(0,-0.20in) -| (rq1.north);
\coordinate (bus) at ([yshift=-0.20in]records.south);
\draw[->] (rq2.north |- bus)--(rq2.north);
\draw[->] (cases.south)--(rq1.north);
\end{tikzpicture}

\caption{Study organization, adapted from the original pipeline figure.
Case construction and agent execution are retained; the former weighted
scoring panel is replaced by outcome assessment and trajectory analysis.
Reference profiles inform analysis, not the agent's instance-specific input.
The lower row is not an execution sequence or a causal model.}
\label{fig:study}
\end{figure*}

\subsection{Case Selection and Reference Validation}

The historical study selected 30 cases from a collection of 696 candidate
IoT CVEs. The selection notes describe three criteria: firmware
accessibility, sufficient public documentation, and membership in one of
three classes. The intended balance is ten buffer overflows, ten command
injections, and ten authentication-bypass or access-control cases. The
case identifiers span 2020--2026. This is a purposive sample for studying
documented reproduction tasks, not a probability sample of IoT
vulnerabilities.

The identifiers match the historical appendix and current archive.
The original case descriptions are retained in
Appendix~\ref{app:case-inventory}, including vendor, product, device class,
and CWE. These are historical metadata awaiting TODO DES-01, not
independently verified facts. The third stratum includes access-control
cases beyond authentication bypass.

Selection-time accessibility and observed acquisition differ: endpoints
change, mirrors contain different versions, and agents must discover URLs.
An acquisition failure does not establish that no affected image exists.

For each CVE, the archive includes a reference profile containing product,
version, vulnerability description, public sources, acquisition hints,
and expected evidence. These profiles were assembled with automated
assistance. They are documentary references, not proof that the authors
have independently replayed every target. A profile may omit a component
or disagree with a later artifact. We retain source provenance and resolve
such disagreements explicitly.

\todo{DES-01}{Reconcile the actual 30 run identifiers with all selection
lists; recover the 696-to-30 exclusion counts and selection dates. Audit
product, affected version, class, ISA, and reference status for every case.}
The completed case table will separate publicly described facts, facts
verified from an acquired image, and independently replayed behavior.
Unknown firmware sizes, dependency requirements, or architectures will
remain unknown rather than being inferred from product names.

\subsection{Agent Configurations and Execution Protocol}

The archive contains five model labels: \texttt{claude-sonnet-4-6},
\texttt{deepseek-v4-flash-free}, \texttt{glm-5.2}, \texttt{gpt-5.5}, and
\texttt{mimo-v2.5-free}. We retain these labels as recorded. In particular,
a route containing \texttt{free} is not assumed to be identical to another
deployment with a similar model name. Each configuration has 90 run
positions, covering 30 CVEs and three repetitions. A shared OpenCode
harness provides the interaction and tool-execution loop~\cite{opencode}.

The historical setup describes an ARM64 host with an NVIDIA GB10 processor
and 128 GB unified memory, and one Ubuntu 24.04 Docker container per run.
Artifacts and traces were collected through separate output volumes.
Independent containers were intended to limit state carryover. The
available setup description identifies the operating system and tool
families but does not fully pin every image, package, and harness version.
\todo{DES-02}{Recover host and container configuration, image digests,
OpenCode commit, tool versions, isolation flags, model providers/snapshots,
sampling parameters, and run dates from original records.}

The common task required an initial plan, execution of the reproduction,
and a final report. It asked agents to record roadblocks, keep outputs in
the designated directory, and stay within a stated \$5 cost and 24-hour
time budget. Agents could install tools and consult public resources.
The six analytical activities were not supplied as a scored checklist.
Removing plan scoring in this study does not remove the planning
instruction from the historical experiments. The unchanged prompt figure
is retained in Appendix~\ref{app:historical-prompt}; its wording does not
explicitly exclude simulation, which matters when interpreting RQ3.

The prompt's stated limits and the orchestrator's enforced limits must be
distinguished. An agent's assertion that a task used only local tools does
not measure model inference cost. Likewise, a timestamp in an agent-written
report need not be an execution timestamp. We use harness metadata where
available and reconcile it with provider and container records.
\todo{DES-03}{Recover the exact executable prompt and control loop;
verify budget enforcement, retry and stop rules, tool permissions,
cross-run reset behavior, and whether any operator interventions occurred.}

\subsection{Evidence Collection and Outcome Assessment}

We index each run by CVE, recorded model label, and repetition. An initial
inventory finds 450 evaluator JSON files, 395 nonempty trace directories,
and 427 nonempty workspace directories. The 55 empty trace directories and
23 empty workspaces are availability observations, not assigned failure
causes. A nonempty directory can still omit the log needed for a specific
judgment. Table~\ref{tab:inventory} records this initial coverage.

\begin{table}[t]
\centering
\begin{tabular}{lrr}
\toprule
Archive object & Present/nonempty & Empty \\
\midrule
Evaluator JSON files & 450 & --- \\
Trace directories & 395 & 55 \\
Workspace directories & 427 & 23 \\
Reference profiles & 30 & --- \\
\bottomrule
\end{tabular}
\caption{Initial file inventory. Nonempty is not a completeness or
correctness judgment. The two directory rows overlap at the run level.}
\label{tab:inventory}
\end{table}

We treat historical evaluator fields as leads for inspection. They do not
replace independent outcome assessment. Positive values on two legacy
triggering items identify 16 candidate runs for review, but the numerical
values are not used as a reproduction threshold. Nor does absence from
that candidate set establish that a run failed: a scorer can omit or
misinterpret evidence. \todo{AUD-01}{Build a run-level evidence manifest,
classify missingness, inspect all claimed successes, and sample other
strata to detect omitted successes and incorrect negative labels.}

Outcome assessment separates three questions. Is the tested artifact an
affected original target? Did the observed input reach it? Does the
observed effect establish the selected vulnerability under the stated
preconditions? A firmware filename alone does not establish the first;
a service banner alone does not establish the second; a script's success
message alone does not establish the third. Evidence should connect
acquisition, target identity, the running process, test input, and effect
within the same run.

\begin{table*}[t]
\centering
\begin{tabular}{p{1.40in}p{2.40in}p{2.50in}}
\toprule
Activity & Retained checks from the original checklist & Evidence-state interpretation \\
\midrule
Information gathering & Product, version, endpoint, vulnerability class,
public sources & Record checked facts and source disagreement, not plan coverage. \\
Firmware acquisition & Image bytes, version metadata, acquisition source
& Distinguish a matching image, a wrong artifact, a failed request, and unknown provenance. \\
Firmware extraction & Container/compression format, extraction output,
usable contents & An extraction command or empty directory does not establish a usable filesystem. \\
Component analysis & ISA/ABI, target binary, vulnerable behavior,
runtime dependencies & Record what was actually identified and the evidence linking it to the CVE. \\
Service rehosting & Runtime, original process, dependencies,
request delivery & Installation or a listening port alone does not prove original-target execution. \\
Vulnerability triggering & Test input, delivery to original target,
specific effect, reusable logs & Apply class-specific checks in
Table~\ref{tab:trigger-evidence}; retain entry-path scope. \\
\bottomrule
\end{tabular}
\caption{Six activities and evidence checks adapted from the original
phase and item tables. Point values, weights, ordered completion, and
attempt bonuses are removed. Each check records supported, contradicted,
unknown, or inapplicable evidence rather than a numerical grade.}
\label{tab:evidence}
\end{table*}

Table~\ref{tab:evidence} retains the original checklist's artifact
questions without its point allocations. Table~\ref{tab:trigger-evidence}
adapts the original trigger-evidence table to the three sampled classes,
adding specificity checks. These are proposed assessment rules; their
comprehensive application and available controls remain pending.

\begin{table*}[t]
\centering
\begin{tabular}{p{1.10in}p{2.45in}p{2.75in}}
\toprule
Selected class & Candidate signal retained from the original table &
Required interpretation check \\
\midrule
Buffer overflow & Fault, core dump, sanitizer report, emulator fault &
Connect the fault to the affected input path; exclude loader or startup
failure. A generic segmentation fault is insufficient. \\
Command injection & Command output or a controlled file side effect &
Attribute the effect to input handled by the original target, excluding
direct creation by the client or setup code. \\
Authentication / access control & Protected operation succeeds under
unauthorized conditions; comparison request & Verify the protected operation,
session or network prerequisites, and a meaningful authorized/denied comparison. \\
\bottomrule
\end{tabular}
\caption{Trigger-evidence criteria adapted from the original appendix,
limited to the studied classes. Benign-input and patched-target controls
can strengthen attribution where feasible; completed controls must be
reported rather than assumed.}
\label{tab:trigger-evidence}
\end{table*}

\todo{AUD-02}{Freeze and apply the per-class rules, revalidate original
target identity and effects, and adjudicate component-level versus
documented-entry reproduction. Report control availability and unresolved
cases. Do not claim that all controls have already been run.}
The final labels will distinguish verified triggering, observed
noncompletion, and insufficient evidence, with scope attached to positive
results. Evidence of simulation and real execution is recorded separately.
No outcome combines different repetitions into a fictitious complete run.

\subsection{Trace Coding and Analysis Strategy}

For RQ2 and RQ3, the intended coding unit is an episode comprising an
action, observable feedback, and subsequent action. A run can contain
several episodes and several obstacle types. We distinguish a visible
obstacle from an analyst's inferred explanation and from an agent's own
statement. For example, a missing-library error is direct execution
feedback; a claim that firmware is unavailable without a recorded search
is an agent assertion with a different evidential status.

An initial codebook uses the six activities and the historical failure
notes as sensitizing concepts. Reviewers will refine the categories on a
pilot sample, preserve ambiguous cases, and record changes before applying
the codebook more broadly. Responses are coded separately from obstacles:
repair, new information gathering, unchanged repetition, target-preserving
alternative, simulation, return to the original target, and stopping.
\todo{COD-01}{Conduct pilot coding; specify analysts, sampling, independent
overlap, agreement measures, codebook revisions, and disagreement
resolution. This draft does not claim completed multi-analyst validation.}

We report both event-level counts and the number of distinct runs
containing each pattern. For a conditional response proportion, the
denominator comprises observable obstacle episodes in the relevant
category; for prevalence across runs, it comprises runs with adequate
evidence for that judgment. This prevents repeated retries in one long
trajectory from being interpreted as many independent affected cases.

RQ1 will report per-run outcomes and the number of CVE--configuration
groups with at least one verified result across three trials. Repetition
is a limited consistency check, not a precise estimate of a case's success
probability. Comparisons should preserve the pairing of configurations
within CVEs; uncertainty estimates must respect CVE clustering.
\todo{ANA-01}{Generate outcome and activity-state tables from adjudicated
records, document the chosen paired/clustered uncertainty procedure, and
report sensitivity to unknown outcomes and incomplete traces.}

For later activities, we will explicitly state which runs reached the
relevant prerequisite. Their conditional results describe that observed
subset. Since different agents can reach different cases, this subset
cannot isolate downstream capability or the causal effect of environment
readiness. Cost and duration likewise require coverage and route-specific
accounting before comparison. \todo{ANA-02}{Audit token, billing, timing,
and termination records; distinguish zero-priced routes from zero
inference use and missing logs from infrastructure failures.}
